\documentclass[10pt]{article}
\usepackage[margin=.8in]{geometry}
\usepackage[T1]{fontenc}
\usepackage{amsmath,amssymb,booktabs,longtable,hyperref,xurl}
\hypersetup{colorlinks=true,urlcolor=blue,linkcolor=blue}
\title{Soft Input Continuity in the Single PCBF--CLF Controller}
\author{Collision Avoidance Research}
\date{October 3, 2026}
\begin{document}
\maketitle
\section{Decision and algorithm}
The user retained the existing 0.10-m optimization margin and requested a soft
preference for gentler maneuvers whenever feasible. The original squared
correction cost favors staying near the initialization inputs, which can
already contain abrupt steering and braking changes. It does not directly
prefer continuity between actual consecutive commands.

The same CLF optimization now minimizes
\[
 \frac{\rho}{\max(1,V(x))}+\epsilon R,\qquad
 R=\frac{R_a+wR_s}{1+w},\quad w=0.5,\quad\epsilon=10^{-4}.
\]
Let $u_i=\bar u_i+d_i$, with componentwise numerical correction radii $r_l$,
$i=0,\ldots,M-1$, $l=1,2$. The anchor term remains
\[
 R_a=\frac{1}{2M}\sum_{i,l}\frac{d_{l,i}^2}{r_l^2}.
\]
For the added continuity preference, set $u_{-1}=\bar u_{-1}$ to the previously
applied command and $d_{-1}=0$. Define
\[
 b_{l,i}=|\bar u_{l,i}-\bar u_{l,i-1}|+r_l+r_{l,i-1},\quad
 B_l^2=\sum_i b_{l,i}^2,\qquad
 R_s=\frac12\sum_l\frac{\sum_i(u_{l,i}-u_{l,i-1})^2}{B_l^2},
\]
where $r_{l,-1}=0$ and the other radii equal $r_l$. Uniform normalization per
actuator preserves equal weighting at every prediction step. The correction
box gives $0\le R_a,R_s,R\le1$, so the existing bound on the effect of the
regularizer on minimum scaled CLF slack remains $\epsilon$ at an exact optimum.
Numerical solver tolerances add their own error. The PCBF budget remains hard.

The anchor reference is still the linearization input, not zero. The new term
vanishes for a constant nonzero sequence matching input memory. There is no
hard steering magnitude or slew constraint, new CLF, controller mode, backup,
cone, optimization stage or relinearization introduced by this change.
Normal frames continue to shift the preceding trajectory; the same existing
failure-triggered initialization and damping remain. The objective is one
sparse positive-semidefinite quadratic with a linear term; its constant is
omitted. Clarabel receives the upper triangular Hessian, as required by its
native adapter. The continuity term couples adjacent inputs. The continuation
state version is 69 and its context includes the new weight.

\section{Validation contract and comparison}
The comparison reuses the original fourteen collision-threat cases, seven
geometries at 8 and 15 m/s, and the original result source commit
\path{b16771139356ec252ed8deb4bb341cc22ffb023a}. Initial ego states, targets and
all pre-existing configuration fields were verified identical. Target
acceleration and sideslip are fixed and known; ego observation is exact.
Road constraints remain absent. Holds are 50 ms, with the same 8/16-step
PCBF prefix and 60-step completion. ODE45 applies the selected command with
relative/absolute tolerances $10^{-11}$ and $10^{-12}$; 31 outputs per hold
independently audit rectangle clearance. Measured computation delay is not
injected into plant time. There is no random seed.

Simulation continues through the post-conflict one-second nominal dwell, with
error tolerances
\[ [0.1\,\mathrm m,1^\circ,0.1\,\mathrm{m/s},
0.05\,\mathrm{m/s},0.01\,\mathrm{rad/s}]. \]
The 1,500-hold observation cap is
not a recovery-speed requirement. The 15-m/s circular head-on fixture has
another encounter near 56 s; that later threat is retained.

For each actuator the measured total variation energy is
\[
 E_l=\frac{1}{h}\sum_{k=0}^{K-1}(u_{l,k}-u_{l,k-1})^2.
\]
The initial input memory is $[0,0]^\top$ in these fixtures. All executed holds
through recovery are included, without averaging by trajectory duration. Thus
longer quiet recovery cannot reduce this metric by diluting a time average.
Body acceleration peaks are evaluated directly from the nonlinear derivative
at the stored dense plant states:
$a_x=\dot v_x-v_y r$, $a_y=\dot v_y+v_x r$. These include one-sided hold
endpoints and are sampled nominal-model measurements.

\section{Results}
All 14 cases recovered without sampled collision or solver interruption over
5,154 holds. Total steering variation decreased by 19.3\% across the family,
and total braking-ratio variation decreased by 38.8\%. Steering variation
improved in 11/14 cases and braking variation in 12/14. Their median per-case
changes were $-6.5\%$ and $-36.8\%$. Peak steering decreased in 12 cases and
peak lateral acceleration in 10 cases. This is a measured preference, not a
claim that every maneuver metric improves.

\begin{center}\small
\begin{tabular}{rlrrrrr}
\toprule
Speed & Scenario & Recovery (s) & Gap (m) & $\Delta E_\delta$ (\%) & $\Delta E_b$ (\%) & Run max (ms)\\
\midrule
8 & headOn & 22.40 & 1.7130 & -54.7 & -64.4 & 78.12 \\
8 & acceleratingHeadOn & 12.70 & 1.3191 & -2.7 & -38.0 & 47.94 \\
8 & brakingLead & 14.85 & 1.3361 & -0.7 & -54.2 & 58.65 \\
8 & crossing & 18.40 & 0.0380 & -42.7 & -79.7 & 39.64 \\
8 & turningCrossing & 12.05 & 0.1039 & -5.4 & -14.3 & 50.93 \\
8 & curvedHeadOn & 27.00 & 2.1486 & -50.3 & -71.5 & 41.19 \\
8 & curvedCrossing & 13.30 & 0.0912 & -5.6 & +4.7 & 53.26 \\
15 & headOn & 10.20 & 0.1791 & -9.6 & -40.9 & 43.05 \\
15 & acceleratingHeadOn & 10.30 & 0.0495 & -7.3 & -45.6 & 74.89 \\
15 & brakingLead & 16.35 & 0.6414 & +12.4 & -23.8 & 59.58 \\
15 & crossing & 11.85 & 3.4016 & +1.2 & -32.5 & 51.58 \\
15 & turningCrossing & 11.55 & 3.1547 & +7.3 & -3.6 & 46.91 \\
15 & curvedHeadOn & 65.95 & 3.9834 & -19.1 & -35.6 & 59.16 \\
15 & curvedCrossing & 10.80 & 2.6241 & -40.0 & +9.5 & 50.24 \\
\bottomrule
\end{tabular}
\end{center}
The 8-m/s head-on case changes from 30.48 to 23.53 degrees maximum steering
and from 7.578 to 5.786 m/s$^2$ peak lateral acceleration. Its recovery
confirmation changes from 12.60 to 22.40 s, which remains within the requested
eventual-recovery goal. The 8-m/s curved head-on case also uses a gentler,
longer maneuver (17.90 to 27.00 s recovery). Conversely, steering variation
increases 12.4\% in the 15-m/s braking-lead case, and its lateral acceleration
peak rises from 7.375 to 7.575 m/s$^2$. Abrupt first avoidance commands can
remain; median change in the maximum sampled steering rate is approximately
zero. No global comfort optimum is claimed.

The 0.10-m setting remains a constraint of the affine collision model. The
smallest dense nonlinear replay gap is 38.044 mm in the 8-m/s crossing case,
down from 49.477 mm as the preceding campaign's overall minimum. In that
particular crossing case the prior gap was about 91.5 mm. A larger nominal
margin has not been silently retained, and this experiment does not establish
a 0.10-m continuous-plant clearance guarantee. Positive measured clearance
is reported separately from the configured margin.

Continuation controller-call median, 95th percentile and maximum are
13.892, 26.420 and 78.116 ms. The prior values were 13.572, 26.399 and
81.396 ms. There are no continuation calls over 100 ms; 24 exceed the 50-ms
sample period. The maximum first call is 887.342 ms, including startup and
cached-construction effects; it is excluded only from the continuation
statistics, not hidden as an achieved real-time startup deadline. Timing was
measured without a profiler on a shared workstation, so this is an observed
distribution rather than a worst-case deadline guarantee.

There are 5,082 one-model frames and 72 failed-shift/fresh-seed frames. A total
of 5,030 frames use one numerical solve, 36 use damping, and the maximum is
four numerical solves in an exceptional frame. Although the preference adds
no optimization stage, it can change active constraints and exceptional solve
counts. The slowest continuation is 8-m/s head-on frame 5:
\begin{center}\begin{tabular}{lr}\toprule Component & ms\\\midrule
Initialization/shift & 1.786\\
Problem construction including CLF & 34.156\\
Numerical optimization & 13.014\\
Full-step model agreement & 0.906\\
Damping & 0.000\\
Other controller overhead & 28.254\\
\bottomrule\end{tabular}\end{center}
The final component is a timing remainder, not a profiled attribution.

\section{Exploration and regression}
Before the user retained 0.10 m, the original objective with a 0.30-m margin
completed recovery in 11 of 14 cases. The 8-m/s curved crossing and both
15-m/s straight head-on variants encountered model-agreement failure.
Short screens at 0.20/0.25 m and a 0.30-m screen with larger initialization
padding did not remove every failure. These observations do not prove physical
infeasibility. None of these margin or padding changes remains in the final
controller.

An initial stage-dependent normalization with weight 9 was discarded. With
uniform channel normalization, weights 0.1, 0.3, 1 and 3 passed short
120-hold screens on three difficult fixtures, but these screens did not prove
recovery. The subsequent full weight-1 campaign recovered in 13/14 cases,
with a model-agreement interruption after 116 holds in the 8-m/s braking-lead
case. Weights 0.1, 0.3 and 0.5 then completed that case; the selected weight
0.5 completed the final full fourteen-case campaign. Selection is empirical
on this scenario family, not a claim of optimal tuning or out-of-family safety.
Exact exploratory outcomes and source hashes are in the adjacent JSON export.

All 395 focused MATLAB tests pass across 18 suites. Added tests verify that
input continuity can improve without sacrificing zero CLF slack or requiring
another solve, that zero weight is permitted and invalid weights are rejected,
and that the solver handles the coupled quadratic's upper-triangular Hessian.
Existing coverage retains shifted budgets, damping, single-CLF behavior,
nonlinear dynamics, terminal continuation and displaced-state recovery.
Code Analyzer reports the pre-existing sparse-index suggestions and an unused
suppression, plus a workstation settings-file warning; no new changed-code
finding was introduced. The source diff passes whitespace validation.

Reproduction uses \texttt{runTimedFlowCampaign(root,out,1500)} with the committed
configuration, the eighteen suites listed in the provenance export, and the
recorded comparison/acceleration scripts. Machine-readable results, tests,
source identities and actual commands accompany this report. Raw experiment
exports and analysis methods are retained as identified archive copies.
\end{document}
